\chapter{Los números racionales}

\begin{objetivos}
  \item Construir los números racionales a partir de los enteros.
  \item Interpretar las fracciones como clases de equivalencia.
  \item Estudiar operaciones, orden y densidad de los racionales.
\end{objetivos}

\section{Motivación: división y fracciones}

En el capítulo anterior se estudió el conjunto de los números enteros, que denotamos por \(\mathbb{Z}\). Este conjunto permite sumar, restar y multiplicar sin salir de él. Sin embargo, no siempre permite dividir.

Por ejemplo, la ecuación
\[
  2x=1
\]
no tiene solución en \(\mathbb{Z}\), porque no existe ningún entero \(x\) cuyo doble sea igual a \(1\). Para resolver ecuaciones de este tipo es necesario ampliar el sistema numérico.

La idea fundamental consiste en introducir objetos que representen cocientes de enteros. Estos objetos reciben el nombre de \emph{números racionales}.

\begin{convencion}[Símbolos básicos]
A lo largo de este capítulo, el símbolo \(\mathbb{Z}\) denota el conjunto de los números enteros. El símbolo \(\mathbb{Q}\), que se definirá más adelante, denotará el conjunto de los números racionales. Todo símbolo que aparezca en una definición será introducido antes de su uso.
\end{convencion}

La expresión \(\frac{a}{b}\), donde \(a\) y \(b\) son enteros y \(b\neq 0\), se interpreta inicialmente como una fracción. El entero \(a\) se llama \emph{numerador} y el entero \(b\) se llama \emph{denominador}.

\begin{advertencia}
La expresión \(\frac{a}{b}\) solo tiene sentido como fracción si el denominador \(b\) es distinto de cero. La división por cero no está definida.
\end{advertencia}

\section{Fracciones y clases de equivalencia}

Una dificultad esencial es que una misma cantidad puede representarse mediante fracciones diferentes. Por ejemplo,
\[
  \frac{1}{2},\qquad \frac{2}{4},\qquad \frac{3}{6}
\]
representan el mismo número racional.

Por tanto, no conviene definir un número racional como una fracción concreta, sino como una clase de fracciones equivalentes.

\begin{definicion}[Conjunto de pares admisibles]
Denotamos por
\[
  S=\{(a,b)\in \mathbb{Z}\times \mathbb{Z}: b\neq 0\}
\]
el conjunto de todos los pares ordenados de enteros cuyo segundo componente es no nulo.
\end{definicion}

Un par \((a,b)\in S\) representa informalmente la fracción \(\frac{a}{b}\).

\begin{definicion}[Equivalencia de fracciones]
Sean \((a,b),(c,d)\in S\). Decimos que \((a,b)\) es equivalente a \((c,d)\), y escribimos
\[
  (a,b)\sim(c,d),
\]
si
\[
  ad=bc.
\]
\end{definicion}

La condición \(ad=bc\) expresa la igualdad usual de fracciones:
\[
  \frac{a}{b}=\frac{c}{d}
  \quad \Longleftrightarrow \quad
  ad=bc.
\]

\begin{proposicion}
La relación \(\sim\) definida sobre \(S\) es una relación de equivalencia.
\end{proposicion}

\begin{proof}
Debemos comprobar que \(\sim\) es reflexiva, simétrica y transitiva.

Sea \((a,b)\in S\). Como \(ab=ba\), se tiene \((a,b)\sim(a,b)\). Por tanto, \(\sim\) es reflexiva.

Supongamos ahora que \((a,b)\sim(c,d)\). Entonces \(ad=bc\). La misma igualdad puede escribirse como \(cb=da\), luego \((c,d)\sim(a,b)\). Por tanto, \(\sim\) es simétrica.

Finalmente, supongamos que \((a,b)\sim(c,d)\) y \((c,d)\sim(e,f)\). Entonces
\[
  ad=bc,
  \qquad
  cf=de.
\]
Multiplicando la primera igualdad por \(f\) y la segunda por \(b\), obtenemos
\[
  adf=bcf,
  \qquad
  bcf=bde.
\]
Por tanto, \(adf=bde\). Como \(d\neq 0\), podemos cancelar \(d\) en \(\mathbb{Z}\) y obtenemos
\[
  af=be.
\]
Así, \((a,b)\sim(e,f)\). Por tanto, \(\sim\) es transitiva.
\end{proof}

\begin{definicion}[Clase de una fracción]
Sea \((a,b)\in S\). La clase de equivalencia de \((a,b)\) es el conjunto
\[
  [(a,b)]=\{(c,d)\in S:(c,d)\sim(a,b)\}.
\]
Esta clase contiene todas las fracciones equivalentes a \(\frac{a}{b}\).
\end{definicion}

Por simplicidad de notación, la clase \([(a,b)]\) se suele escribir como
\[
  \frac{a}{b}.
\]
Esta notación debe entenderse como una clase de equivalencia, no como un único par ordenado.

\begin{ejemplo}
Las fracciones \(\frac{1}{2}\), \(\frac{2}{4}\) y \(\frac{-3}{-6}\) representan el mismo número racional, porque
\[
  1\cdot 4=2\cdot 2,
  \qquad
  1\cdot(-6)=2\cdot(-3).
\]
Por tanto,
\[
  \frac{1}{2}=\frac{2}{4}=\frac{-3}{-6}.
\]
\end{ejemplo}

\section{El conjunto de los números racionales}

Ya podemos definir el conjunto de los números racionales.

\begin{definicion}[Números racionales]
El conjunto de los números racionales, denotado por \(\mathbb{Q}\), es el conjunto de clases de equivalencia de pares \((a,b)\in S\) respecto de la relación \(\sim\). Es decir,
\[
  \mathbb{Q}=S/{\sim}.
\]
Sus elementos se escriben en la forma
\[
  \frac{a}{b},
  \qquad a,b\in\mathbb{Z},\quad b\neq 0.
\]
\end{definicion}

Todo número entero puede verse como un número racional. En efecto, si \(n\in\mathbb{Z}\), entonces
\[
  n=\frac{n}{1}.
\]
De este modo, podemos considerar \(\mathbb{Z}\) como subconjunto de \(\mathbb{Q}\).

\begin{convencion}[Enteros dentro de los racionales]
Cuando escribamos \(n\in\mathbb{Z}\) dentro de un contexto racional, identificaremos \(n\) con el racional \(\frac{n}{1}\).
\end{convencion}

\section{Operaciones con racionales}

Para que \(\mathbb{Q}\) sea un sistema numérico útil, debemos definir suma y producto.

\begin{definicion}[Suma de racionales]
Sean \(\frac{a}{b},\frac{c}{d}\in\mathbb{Q}\), con \(b\neq 0\) y \(d\neq 0\). Definimos su suma por
\[
  \frac{a}{b}+\frac{c}{d}
  =
  \frac{ad+bc}{bd}.
\]
\end{definicion}

\begin{definicion}[Producto de racionales]
Sean \(\frac{a}{b},\frac{c}{d}\in\mathbb{Q}\), con \(b\neq 0\) y \(d\neq 0\). Definimos su producto por
\[
  \frac{a}{b}\cdot\frac{c}{d}
  =
  \frac{ac}{bd}.
\]
\end{definicion}

Estas operaciones están bien definidas: si se cambia una fracción por otra equivalente, el resultado obtenido representa la misma clase racional.

\begin{ejemplo}
Calculamos
\[
  \frac{2}{3}+\frac{5}{4}
  =
  \frac{2\cdot 4+3\cdot 5}{3\cdot 4}
  =
  \frac{8+15}{12}
  =
  \frac{23}{12}.
\]
También,
\[
  \frac{2}{3}\cdot\frac{5}{4}
  =
  \frac{10}{12}
  =
  \frac{5}{6}.
\]
\end{ejemplo}

\begin{proposicion}
El conjunto \(\mathbb{Q}\), con la suma y el producto anteriores, es un cuerpo.
\end{proposicion}

\begin{proof}
La verificación completa consiste en comprobar las propiedades asociativa y conmutativa de la suma y del producto, la existencia de neutro aditivo \(0=\frac{0}{1}\), la existencia de neutro multiplicativo \(1=\frac{1}{1}\), la existencia de opuesto para cada racional y la existencia de inverso multiplicativo para cada racional no nulo. También debe comprobarse la distributividad del producto respecto de la suma.

Estas propiedades se deducen de las propiedades correspondientes de los enteros. En este curso se usará este resultado como una propiedad estructural básica de \(\mathbb{Q}\).
\end{proof}

\begin{advertencia}
La igualdad \(\frac{a}{b}=\frac{c}{d}\) no significa necesariamente que \(a=c\) y \(b=d\). Significa que \(ad=bc\). Por ejemplo, \(\frac{1}{2}=\frac{2}{4}\), aunque \(1\neq 2\) y \(2\neq 4\).
\end{advertencia}

\section{Orden en los racionales}

El orden de los enteros se extiende a los racionales.

\begin{definicion}[Racional positivo]
Sea \(\frac{a}{b}\in\mathbb{Q}\), con \(b\neq 0\). Decimos que \(\frac{a}{b}\) es positivo si \(a\) y \(b\) tienen el mismo signo, es decir, si \(ab>0\).
\end{definicion}

\begin{definicion}[Orden en \(\mathbb{Q}\)]
Sean \(r,s\in\mathbb{Q}\). Decimos que \(r<s\) si \(s-r\) es positivo. Decimos que \(r\leq s\) si \(r<s\) o \(r=s\).
\end{definicion}

En términos de fracciones, si \(b>0\) y \(d>0\), entonces
\[
  \frac{a}{b}<\frac{c}{d}
  \quad \Longleftrightarrow \quad
  ad<bc.
\]

\begin{ejemplo}
Comparamos \(\frac{2}{3}\) y \(\frac{3}{4}\). Como los denominadores son positivos, calculamos
\[
  2\cdot 4=8,
  \qquad
  3\cdot 3=9.
\]
Puesto que \(8<9\), se tiene
\[
  \frac{2}{3}<\frac{3}{4}.
\]
\end{ejemplo}

\begin{advertencia}
La regla de comparar productos cruzados debe aplicarse directamente solo cuando los denominadores son positivos. Si aparece un denominador negativo, conviene cambiar primero la fracción por otra equivalente con denominador positivo.
\end{advertencia}

\section{Densidad de los racionales}

Una propiedad esencial de \(\mathbb{Q}\) es que entre dos racionales distintos siempre hay otro racional.

\begin{teorema}[Densidad de \(\mathbb{Q}\)]
Sean \(r,s\in\mathbb{Q}\) tales que \(r<s\). Entonces existe \(q\in\mathbb{Q}\) tal que
\[
  r<q<s.
\]
\end{teorema}

\begin{proof}
Definimos
\[
  q=\frac{r+s}{2}.
\]
Como \(r\) y \(s\) son racionales, también \(r+s\) es racional, y como \(2\in\mathbb{Z}\) con \(2\neq 0\), el número \(q\) es racional.

Además, de \(r<s\) se deduce
\[
  2r<r+s<2s.
\]
Dividiendo entre \(2\), obtenemos
\[
  r<\frac{r+s}{2}<s.
\]
Por tanto, \(q\) es un racional situado estrictamente entre \(r\) y \(s\).
\end{proof}

\begin{corolario}
Entre dos números racionales distintos existen infinitos números racionales.
\end{corolario}

\begin{proof}
Sean \(r,s\in\mathbb{Q}\) con \(r<s\). Por el teorema anterior existe un racional \(q_1\) tal que \(r<q_1<s\). Aplicando de nuevo el teorema a \(r\) y \(q_1\), existe \(q_2\in\mathbb{Q}\) tal que \(r<q_2<q_1\). Repitiendo el proceso se obtienen indefinidamente racionales distintos entre \(r\) y \(s\).
\end{proof}

\section{Expansiones decimales}

Los números racionales también pueden representarse mediante desarrollos decimales.

\begin{definicion}[Expansión decimal]
Una expansión decimal es una escritura de un número mediante potencias de \(10\). Por ejemplo,
\[
  3.25=3+\frac{2}{10}+\frac{5}{100}.
\]
\end{definicion}

Al dividir dos enteros \(a\) y \(b\), con \(b\neq 0\), el algoritmo de la división produce restos. Como solo hay un número finito de posibles restos, o bien el proceso termina, o bien algún resto se repite.

\begin{proposicion}
Todo número racional tiene una expansión decimal finita o periódica.
\end{proposicion}

\begin{proof}
Sea \(\frac{a}{b}\in\mathbb{Q}\), con \(b>0\). Al realizar la división decimal de \(a\) entre \(b\), los restos posibles son
\[
  0,1,2,\ldots,b-1.
\]
Si en algún momento aparece el resto \(0\), la expansión decimal termina. Si no aparece el resto \(0\), como solo hay \(b-1\) restos no nulos posibles, alguno de ellos debe repetirse. Desde ese momento, las cifras decimales se repiten periódicamente.
\end{proof}

\begin{ejemplo}
El racional \(\frac{1}{4}\) tiene expansión decimal finita:
\[
  \frac{1}{4}=0.25.
\]
El racional \(\frac{1}{3}\) tiene expansión decimal periódica:
\[
  \frac{1}{3}=0.333\ldots
\]
\end{ejemplo}

\begin{erroresfrecuentes}
  \item Pensar que una fracción es el par \((a,b)\), en lugar de la clase de todas las fracciones equivalentes a ella.
  \item Olvidar que el denominador de una fracción no puede ser cero.
  \item Comparar fracciones con denominadores negativos sin transformarlas previamente.
  \item Confundir densidad con continuidad: que entre dos racionales haya otro racional no significa que los racionales llenen todos los puntos de la recta real.
  \item Creer que toda expansión decimal infinita representa necesariamente un número racional. Solo las expansiones decimales finitas o periódicas corresponden a racionales.
\end{erroresfrecuentes}

\begin{seminario}
  \item Operaciones con fracciones como clases de equivalencia.
  \item Problemas de orden y densidad.
  \item Comparación de racionales mediante productos cruzados.
  \item Construcción de racionales entre dos racionales dados.
  \item Análisis de expansiones decimales finitas y periódicas.
\end{seminario}

\begin{ejerciciospropuestos}
  \item Comprueba que \(\frac{6}{8}=\frac{3}{4}\) usando la definición de equivalencia de fracciones.
  \item Calcula \(\frac{2}{5}+\frac{3}{7}\) y \(\frac{2}{5}\cdot\frac{3}{7}\).
  \item Ordena de menor a mayor los racionales \(\frac{1}{2}\), \(\frac{3}{5}\), \(\frac{7}{10}\) y \(\frac{2}{3}\).
  \item Encuentra tres números racionales estrictamente comprendidos entre \(\frac{1}{3}\) y \(\frac{1}{2}\).
  \item Demuestra que si \(r,s\in\mathbb{Q}\) y \(r<s\), entonces \(\frac{2r+s}{3}\) es un racional situado entre \(r\) y \(s\).
  \item Decide si las expansiones decimales siguientes corresponden a números racionales: \(0.125\), \(0.272727\ldots\), \(0.10100100010000\ldots\).
\end{ejerciciospropuestos}

\resumen{El capítulo construye los números racionales como clases de equivalencia de fracciones de enteros. Esta construcción permite definir rigurosamente suma, producto y orden en \(\mathbb{Q}\). Los racionales forman un cuerpo ordenado y tienen la propiedad de densidad: entre dos racionales distintos siempre existe otro racional. Además, sus expansiones decimales son exactamente las finitas o periódicas.}
